\documentclass[10pt]{article}

% Idioma y codificación
\usepackage[utf8]{inputenc}
\usepackage[spanish]{babel}

% Márgenes
\usepackage{geometry}
\geometry{top=2.5cm, bottom=2.5cm, left=2.5cm, right=2.5cm}

% Gráficos y Tablas
\usepackage{pdflscape}
\usepackage{graphicx}
\usepackage{amsmath}
\usepackage{pdfpages}
\usepackage{multicol}
\usepackage{float}
\usepackage[table]{xcolor} % UNA SOLA VEZ Y CON [table]
\usepackage{booktabs}
\usepackage{xltabular}
\usepackage{caption}
\usepackage{subcaption}

% Listings para código
\usepackage{listings}
\lstdefinestyle{terminal}{
	backgroundcolor=\color{gray!10},
	basicstyle=\ttfamily\small,
	keywordstyle=\color{blue},
	commentstyle=\color{gray},
	frame=single,
	columns=fullflexible,
	morecomment=[l]{//}
}
\lstset{
	inputencoding=utf8,
	extendedchars=true,
	literate={á}{{\'a}}1 {é}{{\'e}}1 {í}{{\'i}}1 {ó}{{\'o}}1 {ú}{{\'u}}1 {ñ}{{\~n}}1
}

% Hipervínculos
\usepackage[hidelinks]{hyperref}
\hypersetup{
	colorlinks=true,
	linkcolor=blue,
	filecolor=magenta,
	urlcolor=cyan,
}

% Encabezados y pies de página
\usepackage{fancyhdr}
\pagestyle{fancy}
\fancyhf{}
\fancyfoot[R]{\thepage}

% Formato de títulos
\usepackage{titlesec}
\titleformat{\subsection}{\normalfont\large\bfseries}{\thesubsection}{1em}{}
\titleformat{\subsubsection}{\normalfont\normalsize\bfseries}{\thesubsubsection}{1em}{}

% Datos de la portada
\title{\textit{CARD PARTNER}\\ Una web para amantes de los \textit{TCGS}}

\author{
	\textbf{\small ÁLVAREZ FULCONIS, MARCOS PALATINO}\\
	\textbf{\small LÓPEZ PRIMO, ÁLVARO}\\
	\textbf{\small OTERO GARCÍA, SERGIO}\\
	\textbf{\small SANTIESTEBAN MENDO, EDGAR}
}

\date{Curso 25-26}

\begin{document}
	
	\begin{titlepage}
		\centering
		% Asegúrate de que la ruta al PDF de la UVA sea correcta
		%\includegraphics[scale=0.3]{./Recursos/UVA-3881270087.pdf}\\[0.5cm] 
		\rule{\linewidth}{0.2mm}\\[0.4cm]
		{\huge\bfseries CARD PARTNER\par}
		{\large Una web para amantes de los TCGS\par}
		\rule{\linewidth}{0.2mm}\\[1.5cm]
		{\Large\bfseries Servicios y Sistemas Web}\\[0.3cm]
		{\large\bfseries Laboratorio L2 \\ GRUPO 33}\\[1cm]
		{\Large \begin{tabular}{c} ÁLVAREZ FULCONIS, MARCOS \\ LÓPEZ PRIMO, ÁLVARO \\ OTERO GARCÍA, SERGIO \\ SANTIESTEBAN MENDO, EDGAR \end{tabular}}\\[1.5cm]
		{\Large Curso 25-26}
	\end{titlepage}
	
	\pagenumbering{gobble}
	\tableofcontents \newpage
	
	\pagenumbering{arabic} 
	\setcounter{page}{1}
	
	\section{Presentación del equipo}	
		\subsection{Reparto de carga de trabajo}
			Para esta tercera entrega, hemos decidido cambiar la manera de trabajar, nos hemos divido en dos grupos, Sergio y Marcos se han encargado del \textit{frontend} mientras que Edgar y Álvaro se han encargado de la parte del \textit{backend}.\\
			Esta división se ha hecho con dos propositos, optimizar el tiempo de desarrollo,evitar conflictos en el git y permitir que los integrantes trabajen de la manera más cómoda posible.\\
			Después de la segunda entrega, internamente vimos que Marcos y Sergio tenían más facilidad para el html y los estilos css por lo que se ofrecieron en encargarse de la parte más visual para seguir manteniendo una homogeneidad además de implementar las funcionalidades necesarias (los jsp). En el caso de ser necesario aportar conocimiento concreto sobre las cartas Marcos es un jugador medianamente especializado por lo que podría aportar la información necesaria.\\
			Por otra parte, Edgar y Álvaro se encontraban más cómodos en la gestión de la bbdd, los daos,los dtos y \textit{servlet}, por lo que eligieron de manera voluntaria trabajar en este campo de la web; además, como Álvaro es el que tiene más experiencia en este tcg, construir la base de datos, ha sido una tarea muy sencilla.\\
			De esta manera, al estar en dos equipos, nos hemos podido organizar mejor minimizando así los posibles conflictos en el git puesto que al acotar las tareas, ambas personas del equipo podían trabajar sin pisar al otro integrante.
			
			\begin{table}[h]
				\centering
			
				\begin{tabularx}{\textwidth}{llX} 
					\toprule
					\textbf{Nombre} & \textbf{Porcentaje} & \textbf{Responsabilidad principal} \\
					\midrule
					\multicolumn{3}{l}{\textit{\textbf{Equipo Frontend 50\%}}}  \\
					\midrule
					ÁLVAREZ FULCONIS, MARCOS P. & 25\% & Todo lo relacionado con las vistas de cartas y las del apartado social, navegabilidad de las vistas, maquetación  del documento  \\
					OTERO GARCÍA, SERGIO        & 25\% & Todo lo relacionado con  las vistas de \textit{decks} así como del login, register, página inicial de la web y redacción del documento.  \\
					\midrule
					\multicolumn{3}{l}{\textit{\textbf{Equipo Backend \& Gestión 50\%}}} \\
					\midrule
					LOPEZ PRIMO, ÁLVARO         & 25\% &  Creación del esquema de la bbdd, colaborar en el E/R, población de la base de datos y transformación de los JSONs en datos para la misma, scripts para la automatización de MariaDB en local, redacción del documento. \\
					SANTIESTEBAN MENDO, EDGAR   & 25\% & Colaboración en el E/R, creación de los esquemas de la bbdd, daos y dtos, redacción del documento. \\
					\bottomrule
				\end{tabularx}
					\caption{Distribución de Responsabilidades por Equipo y Miembro}
			\end{table}
	
	
	\section{Introducción: descripción del sistema web}
		PROPUESTA Y ASPECTOS GENERALES (9,5/10) - GRUPO 33\\
		CARD PARNTER ofrece una aplicación que cubre una necesidad real al presentar en una misma página las necesidades más destacadas de los jugadores y aficionados a los TCG. La propuesta presentada, aunque ya existen páginas similares, propone el punto original de unificar dichas necesidades de forma centralizada, recogiendo numerosas funcionalidades para cubrir multitud de las tareas más habituales de los aficionados.\\
		
		DESCRIPCION Y CONTEXTO.(8,5/10)\\
		Se ofrece, de forma detallada y exhaustiva, la descripción y el contexto del problema que se pretende abordar y resolver con esta solución. Queda claro el conocimiento existente al respecto, visto el análisis aportado en el documento. Se citan las referencias a las webs más usadas en este ámbito, aunque se echa en falta comentar más a fondo aspectos de fortalezas y, sobretodo, debilidades de cada una y, que es lo que la solución presentada cubre, añade y/o mejora al respecto.\\
		
	\section{Usuarios objetivo}
		USUARIOS (6,5/10)\\
		Se definen dos tipos de usuario a nivel de 'usuarios finales', con los perfiles claramente definidos. No se cita ningún tipo de usuario con capacidad de administración de la aplicación a nivel de 'backend'/mantenimiento. ¿Existe? Se debe concretar, ya que se está diseñando una aplicación que, posiblemente, requerirá mantenimiento y/o actualizaciones, gestión de los contenidos,... ¿Quién lo lleva a cabo?\\
	
	\section{Especificación de requisitos}
		La especificación de requisitos está organizada y numerada en bloques temáticos, y contiene RNF específicos y mensurables.\\
		
		El modelo de dominio es el que está mejor resuelto: la herencia adecuada de Usuario en Jugador y Administrador, una clase asociativa Contiene con un atributo cantidad, enumeraciones particulares del dominio (tipos, colores,...), y una descripción textual que complementa de manera perfecta el diagrama.\\
		
		El mapa del sitio, que documenta cuatro niveles de profundidad con flujos de navegación detallados, excede lo común en esta entrega. La funcionalidad casi completa la abarcan las páginas estáticas: Deck Builder con ventanas emergentes modales para elegir cartas y líder, visualizador de mazo con precio y winrate, visualizador de carta que incluye todas las características del modelo, así como tres apartados editoriales distintos.\\
		
		Las mejoras que se deben incorporar son: agregar requisitos de información y reglas de negocio explícitas (como la composición mínima de un mazo o las limitaciones de copias por carta); crear la página del perfil del usuario, que incluya la gestión de mazos propios, ya que es el único apartado importante en el mapa que carece de una página; clarificar cómo están relacionados decks.html, community-deck.html y meta-decks.html para prevenir redundancias; y llevar a cabo el enlace con páginas específicas por juego desde la sección TCG.\\
		
		Apartado REQUISITOS:\\
		
		RF-1.1 menciona registro "con cuenta de Google" como alternativa, pero esto no aparece en las páginas estáticas ni en el resto del análisis. Si se mantiene como requisito, debería reflejarse en alguna página.\\
		
		RF-5.2 ("Recomendados: Sugerencia de mazos para principiantes") no tiene una página estática propia clara; los mazos recomendados aparecen en el index pero no como sección independiente.\\
	
	\section{Modelo del dominio}
		
		Apartado MODELO DOMINIO:\\
		
		No se modela la fecha de registro del usuario ni la fecha del comentario (solo fecha\_edicion). Para una web social es relevante distinguir cuándo se creó algo de cuándo se editó por última vez.\\
		
		La relación Guarda entre Jugador y Mazo debería especificar si el jugador puede guardar únicamente mazos públicos de la comunidad o también mazos META oficiales. La multiplicidad está definida pero la restricción semántica no\\
		
		No se modela la relación entre Articulo y las categorías (Noticia, Guía, Torneo) como una enumeración dentro de Articulo con categoria : CategoriaArticulo, cuando en el diagrama sí aparece correctamente. La descripción no lo menciona explícitamente.\\
	
	
	\section{Mapa del sitio web}
		Apartado MAPA WEB:\\
		
		El mapa no incluye explícitamente una página de perfil propio (gestión de datos del usuario, cambio de contraseña) como nodo diferenciado; aparece como "Gestión Mi Cuenta" pero sin desarrollar qué contiene.\\
		
		La sección TCG aparece en el mapa como una opción de ampliación futura, pero ya está implementada en las páginas estáticas. Podría haberse integrado mejor en el mapa como sección funcional actual.\\
	
	
	\section{Estructura de la Base de Datos}
		A espera del feedback
	
	\section{Descripción de la funcionalidad implementada y justificación de la elección}
	
	
	\section{Conclusiones: breve resumen del trabajo desarrollado y aspectos a destacar}
	
	
	\section{Declaración de Uso de Inteligencia Artificial}
		Para esta segunda entrega, se ha usado la IA Gemini como herramienta para dudas y mejoras del archivo LaTeX, revisión de los contenidos y evaluación de los mismos y diseño iterativo para mejorar lo ya realizado.\\
		Además, se ha usado \href{https://gemini.google.com/share/7a923a0bfae8}{Gemini} para la generación del logo del proyecto.\\
		\textbf{Ejemplos de prompts realizados:}
		\begin{itemize}
			\item 'Cómo funciona un servlet en java'
			\item 'Que son los DAOs, DTOs y Servlets en java y como es el flujo de trabajo entre ellos'
			\item 'Como hacer una conexión de MariaDB en ApacheNetBeans'
			\item 'Haciendo push, a mis compañeros no les aparece la bbdd ¿es normal?'
			\item 'Cómo puedo automatizar la instalación y creación de una bbdd con MariaDB'
			\item 'Me puedes revisar los \textit{scripts} de creación de la base de datos'
			\item 'Traduceme este \textit{script} de creación de bbdd para Windows en un .bat'
			\item 'Cómo puedo testear que la base de datos se ha rellenado de manera correcta con los \textit{scrpts} sql'
			\item 'Cómo me puedo saltar la política CORS para hacer peticiones a la web de bandai nanco sobre cartas del tcg de One Piece'
			\item 'Cómo puedo cambiar la numeración de una página para sean botones y mejorar la navegabilidad por una lista de cartas'
	        \item 'Ayudame a mejorar el botón de filtros de la vista cartas'
	        \item '¿Cómo puedo añadir un mensaje de aviso que indique que no se puede acceder a esta vista sin estar loggeado?'
	        \item 'No logro que se apliquen los filtros en las cartas, ¿puedes echarme una mano?'
		\end{itemize}
	
	%\begin{itemize}
	%   \item Herramientas utilizadas: hay que especificar qué herramientas de IA se han utilizado.
	%  \item Propósito del uso: explicar cómo se ha utilizado cada herramienta.
	% \item Alcance del uso: indicar en qué partes específicas del trabajo se ha recibido asistencia de una herramienta de IA.
	%\item Prompts utilizados: incluir los prompts, preguntas o instrucciones exactas que se han
	%introducido en la herramienta de IA.
	%\end{itemize}
	
\end{document}